\section{Consumer surplus and profit comparisons}
\label{sec:comparisons}

The equilibrium correction changes the domain on which the two pricing regimes can be compared, but it does not by itself determine the sign of consumer surplus or profits. Those objects must be recomputed from the realized allocation.

\subsection{Market shares}

At the uniform candidate, the relevant \(A\)-history cutoff is
\[
x_u=\frac12+\frac{s}{6}.
\]
Whenever the uniform profile is a pure equilibrium, Proposition~\ref{prop:uniform} implies \(x_0\ge x_H(s)>x_u(s)\). Thus only \(A\)-history consumers in \((x_u,x_0]\) switch to \(B\), and firm \(A\)'s equilibrium share is \(x_u\).

On the accepted weak-HBP branch, firm \(A\)'s share is \((2-x_0)/3\). Hence, whenever the two pure equilibria coexist,
\begin{equation}
x_u-\frac{2-x_0}{3}
=
\frac{2x_0+s-1}{6}>0.
\label{eq:share-gap}
\end{equation}
The qualitative market-share ranking in accepted Result~2 therefore survives, although the weak-branch displayed difference in the accepted manuscript does not match its own equilibrium shares.

\subsection{Weak-dominance consumer surplus}

The uniform profile has two possible allocation branches even before one asks whether it is an equilibrium. If \(x_0\ge x_u\), some \(A\)-history consumers switch to \(B\). If \(x_0<x_u\), no consumer switches and the inherited allocation remains unchanged. Integrating realized utility over the clipped allocation gives the following profile comparison.

\begin{proposition}[Branch-correct weak-HBP consumer surplus]
\label{prop:weak-cs}
For the accepted weak-HBP price profile and the accepted uniform-price profile,
\[
CS^d-CS^u=
\begin{cases}
\displaystyle
\frac{-\tau^2(52x_0^2-52x_0-1)
+2\sigma\tau(18x_0-17)+\sigma^2}{36\tau},
&x_0\ge x_u,\\[1.2em]
\displaystyle
\frac{\sigma^2+12\sigma\tau x_0-14\sigma\tau
-8\tau^2x_0^2+8\tau^2x_0+5\tau^2}{18\tau},
&x_0<x_u.
\end{cases}
\]
The two expressions agree at \(x_0=x_u\), and \(CS^d-CS^u>0\) throughout the accepted weak-dominance profile domain \(1/2<x_0<\bar x(s)\).
\end{proposition}

The first branch differs algebraically from accepted Eq.~(15). More importantly, the accepted one-way-switch integral cannot be extended into \(x_0<x_u\), because its switching interval is then reversed. The exact point \((s,x_0)=(0.99,0.501)\), for example, lies on the no-switch branch and gives
\[
CS^d-CS^u=\frac{17993}{4500000}>0.
\]
At a valid switching-branch point inside the corrected pure-equilibrium overlap, \((s,x_0)=(0.1,0.7)\), direct integration yields \(221/720\), while the accepted Eq.~(15) expression evaluates to \(1279/3600\). Both are positive, but they are not equal.

\begin{corollary}[Weak-HBP equilibrium consumer surplus]
\label{cor:weak-cs}
Whenever weak HBP and uniform pricing both admit the pure equilibria being compared,
\[
CS^d>CS^u.
\]
Thus the consumer-surplus sign reversal stated in accepted Result~3 is not reproduced on the corrected pure-equilibrium domain.
\end{corollary}

\subsection{The smaller firm's profit under weak dominance}

Using the accepted weak-HBP prices and the corrected uniform equilibrium prices, the normalized profit difference for firm \(B\) is
\begin{equation}
\frac{\pi_B^d-\pi_B^u}{\tau}
=
\frac{s^2-12sx_0+14s+20x_0^2-20x_0+1}{18}.
\label{eq:B-profit-gap}
\end{equation}

\begin{proposition}[Smaller-firm profit on the pure overlap]
\label{prop:B-profit}
For every \(0\le s<s_c\) and every \(x_H(s)\le x_0<\bar x(s)\),
\[
\pi_B^d<\pi_B^u.
\]
\end{proposition}

The numerator of \eqref{eq:B-profit-gap} is convex in \(x_0\). On the larger interval \([1/2,\bar x]\), its endpoint values are
\[
s^2+8s-4<0
\quad\text{and}\quad
\frac{(1+s)(21s-11)}4<0
\]
for \(0\le s<s_c\). Convexity then implies negativity throughout the interval. Consequently, the high-switching-cost region associated with accepted Result~4 can still be evaluated against the Eq.~(12) counterfactual profile, but it is not a comparison of two pure equilibria.

\subsection{What the equilibrium correction does}

The distinction between profile and equilibrium comparisons is central. Proposition~\ref{prop:weak-cs} is a statement about two specified price profiles on the entire source weak-dominance domain. Corollary~\ref{cor:weak-cs} and Proposition~\ref{prop:B-profit} are equilibrium statements only on
\[
0<s<s_c,
\qquad
x_H(s)\le x_0<\bar x(s).
\]
No consumer-surplus or profit value is attributed to an uncharacterized mixed uniform-pricing equilibrium below \(x_H\).
